\chapter{Tuning Resources and Configuration}
\label{ch:tuning}

\begin{chapterintro}
This chapter does not hand you a table of recommended settings. It derives an
executor shape from the memory map in \Cref{ch:execmem} and the shuffle in
\Cref{ch:shuffle}, so that every number --- cores, heap, overhead, partition
count --- has a reason you can defend. It then covers dynamic allocation, the
external shuffle service on Kubernetes, configuration precedence, and running
Spark behind Kyuubi for many users. For the standalone list of tuning knobs with
values, see \texttt{spark\_memory\_tuning.md} in the repository root.
\end{chapterintro}

\section{Executor sizing from the mechanism}
\label{sec:executor-sizing}
\tbd

\section{Dynamic allocation}
\label{sec:dynamic-allocation}
\tbd

\section{Parallelism defaults}
\tbd

\section{The external and push-based shuffle service}
\tbd

\section{Configuration precedence and profiles}
\tbd

\section{Kyuubi and multi-tenancy}
\tbd

\section{Summary}
\tbd

\exercisehead
\begin{exercises}
\item Take a job that spills. Propose two different fixes --- more memory versus
      fewer cores --- and predict which the mechanism favours, and why.
\end{exercises}
